% ASCE Journal Article Template
% Suitable for: Journal of Structural Engineering, Journal of Bridge Engineering,
%               Journal of Computing in Civil Engineering, Journal of Engineering Mechanics,
%               Structural Health Monitoring, Journal of Geotechnical Engineering, etc.
% Last updated: 2024
%
% ASCE citation style: author-year, e.g. (Smith and Jones 2020)
% Reference format: Author, F. M. (Year). "Title of paper." Journal Name Volume(Issue), Pages.
%                   https://doi.org/xxxxx

\documentclass[12pt]{article}

% ── Core packages ─────────────────────────────────────────────────────────────
\usepackage[margin=2.54cm]{geometry}   % 1-inch margins (ASCE requirement)
\usepackage{times}                      % Times New Roman
\usepackage{setspace}                   % Line spacing
\usepackage{graphicx}                   % Figures
\usepackage{amsmath,amssymb}            % Math
\usepackage{booktabs}                   % Professional tables
\usepackage{caption}                    % Caption formatting
\usepackage{subcaption}                 % Subfigures
\usepackage{multirow}                   % Multi-row table cells
\usepackage{hyperref}                   % Hyperlinks / DOI
\usepackage{lineno}                     % Line numbers (required for review)
\usepackage{siunitx}                    % SI units (\SI{25}{\mega\pascal})
\usepackage[round,sort,authoryear]{natbib}  % ASCE author-year citations

% ── Line numbering ─────────────────────────────────────────────────────────────
\doublespacing
\linenumbers

% ── Caption style (ASCE: bold label, period separator) ────────────────────────
\captionsetup{labelfont=bf, labelsep=period, font=small}

% ── Hyperref options ──────────────────────────────────────────────────────────
\hypersetup{colorlinks=true, linkcolor=black, citecolor=black, urlcolor=blue}

% ══════════════════════════════════════════════════════════════════════════════
\title{%
  Insert Your Title Here: Concise, Descriptive, and Reflecting Main Contribution
}

\author{
  First Author\textsuperscript{1},
  Second Author\textsuperscript{2}, and
  Corresponding Author\textsuperscript{3}
}

\date{}

\begin{document}
\maketitle

% ── Affiliations ──────────────────────────────────────────────────────────────
\noindent
\textsuperscript{1}Ph.D. Candidate, Dept.\ of Civil and Structural Engineering,
University Name, City, State/Province Postal Code, Country.
E-mail: \href{mailto:first@university.edu}{first@university.edu} \\[2pt]
\textsuperscript{2}Associate Professor, Dept.\ of Civil Engineering,
University Name, City, Country. \\[2pt]
\textsuperscript{3}Professor, Dept.\ of Civil Engineering, University Name,
City, Country (corresponding). E-mail: \href{mailto:corresponding@university.edu}{corresponding@university.edu}

\bigskip

% ── Abstract ─────────────────────────────────────────────────────────────────
% ASCE: ≤250 words, structured or unstructured depending on journal
\begin{abstract}
\noindent
Write a concise abstract (≤250 words) covering: (1) motivation and research gap,
(2) objectives, (3) methods (briefly), (4) main results with quantitative values,
and (5) key conclusions and practical significance.
Avoid references, abbreviations, and jargon in the abstract.
State specific performance metrics (e.g., \SI{15}{\percent} increase in load capacity,
$R^2 = 0.95$ prediction accuracy) rather than vague claims.
\end{abstract}

% ── Keywords ─────────────────────────────────────────────────────────────────
\noindent\textbf{Keywords:} keyword one; keyword two; keyword three;
keyword four; keyword five

\newpage

% ══════════════════════════════════════════════════════════════════════════════
\section{Introduction}
% Recommended length: 600–900 words
% Paragraph structure:
%   (1) Broad context and engineering significance
%   (2) Literature review — what has been done and by whom
%   (3) Research gap — what is missing / what problem remains
%   (4) Objectives — what this paper does
%   (5) Paper organization (optional for longer papers)

Provide engineering context and motivation.
Review relevant prior work, citing key references using author-year style:
\citet{Author2023} demonstrated that \ldots, while \citet{Smith2021} showed \ldots
For parenthetical citations use \citep{Jones2022}.

State the research gap clearly.
Present your specific objectives.

% ══════════════════════════════════════════════════════════════════════════════
\section{Background and Literature Review}
% Optional — merge into Introduction for shorter papers

Summarize the state of knowledge relevant to this study.
Organize thematically (not chronologically).

% ══════════════════════════════════════════════════════════════════════════════
\section{Methodology}
% Alt titles: "Numerical Model", "Experimental Program", "Analytical Framework"

\subsection{Problem Description and Scope}
Describe the structure, material, system, or phenomenon under investigation.
Report geometry, material properties, loading conditions, and boundary conditions.
Use SI units with \texttt{siunitx}: e.g., $f'_c = \SI{30}{\mega\pascal}$,
$f_y = \SI{420}{\mega\pascal}$, story height $= \SI{3.5}{\meter}$.

\subsection{Finite Element Model}
% Remove/rename if not applicable
Describe element types, mesh strategy, constitutive models, and convergence criteria.
Report software and version (e.g., ABAQUS 2023, OpenSees 3.5.0, ANSYS 2024 R1).
Mesh convergence must be demonstrated with at least three mesh densities.
Validate model against experimental data before parametric study.

\subsection{Experimental Setup}
% Remove/rename if not applicable
Describe specimen geometry, instrumentation (gauge locations, sampling rate),
loading protocol (monotonic vs.\ cyclic, loading rate), and test matrix.
Follow ASTM/ASCE/ACI standards for material testing procedures.

\subsection{Validation}
Compare model predictions against experimental or analytical benchmarks.
Report error metrics: mean absolute error (MAE), root-mean-square error (RMSE),
and coefficient of variation (CoV).

% ══════════════════════════════════════════════════════════════════════════════
\section{Results}
% Present results without interpretation; save interpretation for Discussion
% Organize by finding, not by sub-task

\subsection{Result Topic One}
Present quantitative results.
Refer to figures and tables using standard LaTeX references:
see Fig.~\ref{fig:example} and Table~\ref{tab:example}.

\begin{figure}[htbp]
  \centering
  \includegraphics[width=0.75\columnwidth]{figures/figure_01.png}
  \caption{Caption describing the figure. Include axis labels, units, and key observations.
           Panels labeled (a), (b), etc.\ for multi-panel figures.}
  \label{fig:example}
\end{figure}

\begin{table}[htbp]
  \centering
  \caption{Table title. Units in column headers. Significant figures consistent.}
  \label{tab:example}
  \begin{tabular}{lccc}
    \toprule
    Specimen & $f'_c$ (\si{\mega\pascal}) & $V_n$ (\si{\kilo\newton}) & Failure mode \\
    \midrule
    S1       & 30.2 & 245.6 & Flexure \\
    S2       & 40.5 & 312.1 & Shear    \\
    S3       & 50.0 & 380.4 & Flexure  \\
    \bottomrule
  \end{tabular}
  \begin{tablenotes}
    \small
    \item Note: $V_n$ = nominal shear strength; failure mode per ACI 318-19 definitions.
  \end{tablenotes}
\end{table}

\subsection{Result Topic Two}
Continue presenting results.

% ══════════════════════════════════════════════════════════════════════════════
\section{Discussion}
% Interpret results in the context of research objectives and prior literature
% Address: (1) mechanism explanation, (2) comparison with literature,
%           (3) practical implications, (4) limitations

Interpret the findings.
Compare with results from the literature \citep{Author2020}.
Discuss practical implications for engineering design or practice.
Acknowledge limitations of the study.

% ══════════════════════════════════════════════════════════════════════════════
\section{Conclusions}
% Numbered list of conclusions is STRONGLY preferred by ASCE
% Each item should be concise, specific, and quantitative where possible

The following conclusions are drawn based on the results of this study:
\begin{enumerate}
  \item First conclusion — state the main finding with specific numbers.
  \item Second conclusion — relate to practical design implications.
  \item Third conclusion — note limitations or conditions of applicability.
  \item Future work recommendation (optional).
\end{enumerate}

% ══════════════════════════════════════════════════════════════════════════════
% Data Availability Statement (required by ASCE since 2022)
\section*{Data Availability Statement}
All data, models, and code that support the findings of this study are available
from the corresponding author upon reasonable request.
% OR:
% Data are available in the repository: [repository name and DOI/URL].

% ══════════════════════════════════════════════════════════════════════════════
\section*{Acknowledgments}
The authors acknowledge financial support from [funding agency] under
Grant No.\ [grant number]. The authors declare no conflicts of interest.

% ══════════════════════════════════════════════════════════════════════════════
% ASCE REFERENCE FORMAT:
% Author, F. M. (Year). "Title of paper." Journal Name Volume(Issue), Pages.
%   https://doi.org/xxxxx
%
% Example BibTeX entries for ASCE format:
% @article{Smith2021,
%   author  = {Smith, John A. and Jones, Mary B.},
%   title   = {Seismic Performance of RC Shear Walls with Boundary Elements},
%   journal = {J. Struct. Eng.},
%   year    = {2021},
%   volume  = {147},
%   number  = {3},
%   pages   = {04021010},
%   doi     = {10.1061/(ASCE)ST.1943-541X.0002934}
% }

\bibliographystyle{abbrvnat}   % Author-year, abbreviated journal names
\bibliography{../references/references}

\end{document}
